\section{Ethical Considerations}
\label{sec:ethics}

\emph{Deployment context and authorization.} The system operates a live
deception honeypot against real traffic, which obligates care. The
testbed is isolated from production networks and reachable only over a
private overlay (tailnet); the deception targets only actors actively
attacking it, and a red-team authorization context governed the project's
own probing of the stack. All credential material shown to attackers is
registered canary secrets: nothing real is ever exposed, and any canary
appearing in outbound traffic is attributable to bait interaction by
construction. The human-in-the-loop gate keeps a human decision between
the system and the highest-stakes content classes, with every gate
decision recorded in an append-only audit log.

\emph{Committed data contains no real-attacker session logs.} The released
corpora are driver-generated: the 13 replayed real-log sessions derive
from the project's own Cowrie captures of traffic against the isolated
testbed (not production services), contain no production credentials, and
are released only in feature form (command sequences, family labels), not
as raw transcripts. No logs of real external-attacker sessions are
committed to the artifact; everything a reviewer can download was produced
by the benchmark driver.

\emph{Identifier handling.} Cross-session attacker profiles are keyed by
a hash of the source IP computed as HMAC-SHA256 under a deployment-provided
secret key; no default key ships with the repository, so the released
profiles cannot be linked back to originating addresses by anyone who
lacks the operator's secret. The paper reports no IP addresses.

\emph{Assessment is by an LLM-judge panel, not human raters.} Deception
quality is scored by a three-persona panel of independently prompted LLM
judges sharing one judge model (\S\ref{sec:rq3}). We state plainly that
no human raters participated: the $0.750$ agreement figure is a
within-model consistency measure, not inter-annotator agreement, and the
$3.86/5$ consensus is one evaluation instrument's verdict, not a human
usability finding. A small expert-rater pass over the eight released
transcripts is named as required future work.

\emph{Dual use.} The system is an AI-driven deception architecture, and
the same machinery that deceives an attacker into a controlled fiction
could, in other hands, deceive a user into trusting a compromised system;
conversely, the release includes the verifier's decision procedures and
scenario sets, so an attacker who studies them knows exactly which
consistency surface is enforced and can craft probes against the
extraction layer's recall. We judge the reproduction and audit value of
the release to outweigh this exposure, but the surface is now declared.
We release scenario sets and scoring harnesses so the community can audit
the evaluation, not just the code.

\emph{Open fingerprinting channels.} The system's deception is honest
about what it does not hide: response latency (generative turns average
tens of seconds) and cross-session divergence of the fiction remain open
fingerprinting channels by design admission
(\S\ref{sec:threatmodel}), and we treat homogenizing either as future
work rather than a solved problem.
